\documentclass[10pt]{article}

\usepackage[utf8]{inputenc}

\usepackage[T1]{fontenc}

\usepackage{amsmath}

\usepackage{amsfonts}

\usepackage{amssymb}

\usepackage[version=4]{mhchem}

\usepackage{stmaryrd}

\usepackage{mathrsfs}

\usepackage{graphicx}

\usepackage[export]{adjustbox}

\graphicspath{ {./images/} }



\begin{document}

тов $a_{j}(t)$ регулярна для уравнений (1) тогда и только тогда, когда все функции $(t-a)^{j} a_{j}(t), j=1, \ldots, n$, голоморфыы в нуле. Для систем (2) справедливо следующее достаточное условие: если әлементы матрицы $A(t)$ имеют простой полюс в точке $a$, то эта точка - Р. о. т. системы (2). Явное условие на матрицу $A(t)$, необходимое и достаточное для того, чтобы точка $a$ была Р. о.т. системы (2), пока (1983) не получено.



Jит.: [1] $\Gamma$ о л у б е в В. В., Лекции по аналитической теории дифференциальных уравнений, 2 изД., М.- ЈI., 1950; [2] К о ддинг тон Э. А., Ле в и н с о н Н., Теория обыкновенных дифференциальных уравнений, пер. с англ., М., 1958; [3] L еv e l t A. H. M., "Proc. Koninkl. Nederl. akad. wet. Ser. A», 1961, v. 64, № 4, p. 362-403; [4] D e l i g n e P., Equations différentielles à points singuliers réguiliers, B., 1970 (Lect. Notes in Math., № 163); [5]-P l e m e $1 \mathrm{j}$ J., Problems in the sense of Riemann and Klein, Univ. of Adelaide, 1964.



РЕГУЛЯРНАЯ ПОЛУ ГРУППА - полугруппа, каждый элемент к-рой регулярен.



Произвольная Р. п. $S$ содержит идемпотенты (см. Регулярный элемент), и строение $S$ в значительной степени определяется «строением» и "расположением» в $S$ множества всех ее идемпотентов $E(S)$. Р. п. с единственным идемпотентом - это в точности группы. На $E(S)$ прежде всего можно смотреть как на частично упорядоченное множество относительно естественного частичного порядка (см. Идемпотент); известны структурные теоремы, описывающие Р. п. $S$ с нек-рыми естественными ограничениями на множество $E(S)$. Одно из таких ограничений (для полугрупп с нулем) состоит в том, что все ненулевые идемпотенты примитивны (см. Вполне простая полугруппа); полугруппа с этим свойством наз. п р и м и т и в н о й. Следующие условия для полугруппы $S$ әквивалентны: а) $S$ есть примитивная Р. п., б) $S$ есть Р. п., равная объединению своих 0-минимальных (см. Минимальный идеал) правых идеалов, в) $S$ есть 0 -прямое объединение вполне 0-простых полугрупп. Описано строение Р. п., у к-рых $E(S)$ есть цепь, упорядоченная по типу отрицательных целых чисел [2].



Более информативный взгляд на $E(S)$ состоит в рассмотрении на этом множестве частичной операции о, заданной следующим образом. Если для $e, f \in E(S)$ хотя бы одно из произведений $e f, f e$ равно одному из әлементов $e, f$, то $e f \in E(S)$; полагают тогда $e \circ f=e f$. Возникающая частичная алгебра может быть охарактеризована аксиомами, использующими два отношения квазипорядка $\omega^{r}$ и $\omega^{l}$, тесно связанные с заданной частичной операцией (реализация этих отношений в $E(S)$ такова: $e \omega^{r} f$ означает $f e=e, e \omega^{l} f$ означает $e f=e$; тогда $\omega^{r} \cap \omega^{l}$ есть отношение естественного частичного порядка на $E(S))$; такая частичная алгебра наз. б иу по р я до ч е нны м мно же с т во м (см. [5]). Произвольная Р. п. может быть определенным образом сконструирована из биупорядоченного множества и групп. Таким образом, в терминах биупорядоченных множеств можно проводить классификацию Р. п. Среди исследованных в этом направлении типов полугрупп - ком б и на то рны е Р. п. (см. [7]), т. е. имеющие лишь одноэлементные подгруппы.



Гомоморфный образ Р. п. будет Р. п. Всякий нормальный комплекс Р. п., являющийся подполугруппой, содержит идемпотент. Произвольная конгруэнция на Р. п. однозначно определяется своими классами, содержащими идемпотенты. Конгруэнция на Р. п. $S$ разделяет идемпотенты тогда и только тогда, когда она содержится в отношении $\mathscr{H}$ (см. Грина отношения эквивалентности); множество таких конгруәнций составляет модулярную подрешетку с нулем и единицей в решетке всех конгруәнций на $S$. Р. п. наз. ф у н д а м е нт а л ь н о й, если эта подрешетка состоит лишь из отношения равенства. Всякая комбинаторная Р. П. будет фундаментальной Фундаментальные Р. П. важны не только как один из более обозримых типов Р. п., но и в силу определенной «универсальности» их класса для Р. п. А именно, для любого биупорядоченного множества $E$ можно канонич. образом сконструировать фундаментальную Р. п. $T_{E}$, для к-рой $E$ будет биупорядоченным множеством всех идемпотентов, причем для любой Р. п. $S$ такой, что $E(S)=E$, существует разделяющий идемпотенты гомоморфизм $\varphi: S \rightarrow T_{E}$, для к-рого $\varphi(S)$ будет подполугруппой в $T_{E}$, содержащей $E$ (о различных конструкциях для $T_{E}$ см. [3], [5], [8], [10]). P. п. $S$ фундаментальна тогда и только тогда, когда гомоморфизм $\varphi$ инъективен.



Если $S-$ Р. п., то подполугруппа $\langle E(S)\rangle$, порожденная всеми ее идемпотентами, также будет Р. п. Подполугруппа $\langle E(S)\rangle$ оказывает существенное влияние на строение $S$. Р. П. идемпотентно порождена тогда и только тогда, когда таков каждый ее главный фактор [10]. В идемпотентно порожденной Р. п. $S$ для произвольного элемента $x$ существует представление $x=$ $=e_{1} e_{2} \ldots e_{n}$, где $e_{i} \in E(S)$ и $e_{i}(\mathscr{L} \bigcup \mathscr{R}) e_{i+1}$ при $i=$ $=1, . ., n-1$ (здесь $\mathscr{L}$ и $\mathscr{R}$ - отношения эквивалентности Грина) (см. [5]). Последовательность идемпотентов $e_{1}, e_{2}, \ldots, e_{n}$ с указанным свойством наз. $E$ - ц еп ь ю. В бипростой идемпотентно порожденной полугруппе любые два идемпотента связаны нек-рой $E$ цепью, и если они сравнимы в смысле естественного частичного порядка, то длина такой цепи $\geqslant 4$.



Если $\langle E(S)\rangle=E(S)$, т. е. произведение любых двух идемпотентов снова есть идемпотент, то Р. п. $S$ наз. о р т о д о к с а л в н о̆̈̈. Класс ортодоксальных полугрупп содержит, в частности, все инверсные полугруппы. Полугруппа ортодоксальна тогда и только тогда, когда каждый ее главный фактор ортодоксален. Имеются структурные теоремы об ортодоксальных полугруппах (см. [4], [9]).



Отношение естественного частичного порядка на $E(S)$ может быть продолжено на Р. п. $S$ следующим образом: $x \leqslant y$, если существуют идемпотенты $e$ и $f$ такие, что $x=e y=y f$. В случае, когда $S$ инверсна, отношение $\leqslant$ превращается в отношение естественного частичного порядка, для произвольной Р. п. оно также наз. отнонением е с т е с т в е н і о г о ч а с т и чв о г о по р я д к а. Отношение $\leqslant$ на Р. п. $S$ corласовано с умножением тогда и только тогда, когда для любого идемпотента $e$ подполугруппа $e S e$ инверсна [6]. Р. п., обладающие таким свойством, наз. п с е вд о и в в е р с ны ми. Более широкий класс составляют п с е в д о о р т о д о к с а ль н ы е полугруппы (для любого идемпотента $e$ подполугруппа $e S e$ ортодоксальна). Для указанных классов полугрупп использовались также термины «локально инверсная» и «локально ортодоксальная». Р. П. наз. е с т е с т в е нн о й, если множество всех ее групповых әлементов (см. Регулярный элемент) есть подполугруппа. Имеются структурные теоремы о псевдоинверсных, псевдоортодоксальных [11] и естественных [12] Р. п.



Многие структурные теоремы о различных типах Р.п. представляют собой (подчас весьма далекие) обобщения и модификации конструкции рисовской полугруппы матричного типа или суммы прямого спектра групп (см. Клифбордова полугруппа), либо оюираются на те или иные представления полугрупп и разложения их в подпрямые произведения (см. [1], [13]). См. также ст. Il олугруппа.



Лит.: [1] К л и ф фо р д А., П р е с т о н Г., Алгебраическая теория полугрупп, пер. с англ., т. 1-2, М., 1972; [2] M u n n W. D., "Glasgow Math. J.", 1968 , v. 9, pt. 1, p. 46-66; [3] C 1 i fж e, «J. pure and appl. algebra», 1976, v. 8, p. 23-50; [5] N a mb o o r i p a d K. S. S., "Mem. Amer. Math. Soc..", 1979, v ${ }^{22}$, v. 23, pt 3, p. 249-60; [7] "N a m b o o r i p a d K. S. S. S., R a j a n A. R., "Quart. J. Math.", 1978, v. 29, № 116, p. 489-504; [8]



\begin{center}

\includegraphics[max width=\textwidth]{2023_07_08_3792c89db3478732fd5fg-1(1)}

\end{center}



\begin{center}

\includegraphics[max width=\textwidth]{2023_07_08_3792c89db3478732fd5fg-1}

\end{center}





\end{document}